%------------------------------------------------------------------------------%

A Tabela~\ref{tab:series-maclaurin} lista algumas séries de Maclaurin
importantes e seus raios de convergência.

%------------------------------------------------------------------------------%
\begin{ntable}[\setlength{\extrarowheight}{5mm}]%
              {@{\extracolsep{-5pt}}%
               >{$\displaystyle\vphantom{\int_\int^\int}}w{c}{0.12\linewidth}<{$}%
               >{$\displaystyle\vphantom{\int_\int^\int}}w{l}{0.23\linewidth}<{$}%
               >{$\displaystyle\vphantom{\int_\int^\int}}w{l}{0.45\linewidth}<{$}%
               >{$\displaystyle\vphantom{\int_\int^\int}}w{l}{0.10\linewidth}<{$}}[1][5mm]
              {Séries de Maclaurin}
              {tab:series-maclaurin}
  \frac{1}{1-x}
  & = \sum_{n=0}^{\infty} x^n
  & = 1 + x + x^2 + x^3 + \cdots
  & R = 1
  \\ \hline
  e^x
  & = \sum_{n=0}^{\infty} \frac{x^n}{n!}
  & = 1 + \frac{x}{1}
    + \frac{x^2}{2!}
    + \frac{x^3}{3!}
    + \cdots
  & R = \infty
  \\ \hline
  \sin(x)
  & = \sum_{n=0}^{\infty} \frac{(-1)^n x^{2n+1}}{(2n+1)!}
  & = x - \frac{x^3}{3!}
    + \frac{x^5}{5!}
    - \frac{x^7}{7!}
    + \cdots
  & R = \infty
  \\ \hline
  \cos(x)
  & = \sum_{n=0}^{\infty} \frac{(-1)^nx^{2n}}{(2n)!}
  & = 1 - \frac{x^2}{2}
    + \frac{x^4}{4!}
    - \frac{x^6}{6!}
    + \cdots
  & R = \infty
  \\ \hline
  \arctan(x)
  & = \sum_{n=0}^{\infty} \frac{(-1)^n x^{2n+1}}{2n+1}
  & =	x - \frac{x^3}{3}
    + \frac{x^5}{5}
    - \frac{x^7}{7}
    + \cdots
  & R = 1
  \\ \hline
  \ln(x+1)
  & = \sum_{n=1}^{\infty} \frac{(-1)^{n-1}x^n}{n}
  & = x - \frac{x^2}{2}
    + \frac{x^3}{3}
    - \frac{x^4}{4}
    + \cdots
  & R = 1
  \\ \hline
  (1+x)^k
  & = \sum_{n=0}^{\infty} \binom{k}{n} x^n
  & = 1 + kx
    + \frac{k(k-1)}{2}x^2
    %	+ \frac{k(k-1)(k-2)}{3!}
    + \cdots \quad
  & R = 1
\end{ntable}
%------------------------------------------------------------------------------%

\FloatBarrier

%------------------------------------------------------------------------------%
